\documentclass[10pt,a4paper]{amsart}
\usepackage[margin=19mm]{geometry}
\usepackage{amsmath,amssymb,mathtools}
\usepackage[scr=rsfs,cal=euler]{mathalfa}
\usepackage{xfrac}
\usepackage[unicode,hidelinks,bookmarks=false]{hyperref}
\newcommand{\ie}{i.e.\ }
\newcommand{\cf}{cf.\ }
\newcommand{\resp}{respectively\ }
\newcommand{\Subsection}{\subsection}
\providecommand{\nobreakdash}{\nobreak\hbox{-}}
\setlength{\emergencystretch}{2em}
\multlinegap0pt
\usepackage{bm}
\usepackage{bbm}
\usepackage{dsfont}
\usepackage{xspace}
\usepackage{cancel}
\usepackage{mathtools}
\usepackage{amsthm}
\makeatletter
\@ifundefined{theorem}{\newtheorem{theorem}{Theorem}}{}
\@ifundefined{proposition}{\newtheorem{proposition}[theorem]{Proposition}}{}
\makeatother
\newcommand{\arc}{\mathrm{arc}}
\title[Decreasing paths of polygons]{Decreasing paths of polygons}
\author{Isaac Kulp, Charlotte Ochanine, Logan Richard, Leonel Robert, Scott Whitman}
\date{Source: arXiv:2402.12643v2 (CC BY 4.0)}
\begin{document}
\maketitle

\noindent\emph{Source and licence.} Decreasing paths of polygons. Version arXiv:2402.12643v2 (CC BY 4.0), distributed under the Creative Commons Attribution 4.0 International licence, as recorded in the arXiv licence metadata for this exact version and shown on the abstract page https://arxiv.org/abs/2402.12643v2. Full rights notice: licenses/arxiv-2402.12643-CC-BY-4.0.txt.\par\smallskip

\noindent\textit{Editorial scope.} Counted unit: the Proposition whose source label is \texttt{piorientation} in the article (source0.tex lines 1591--1593, source label \texttt{piorientation}), delivered together with the article's own complete original proof of it (source0.tex lines 1595--1648, one \texttt{proof} environment of 54 lines). No mathematics is written, repaired or re-derived: statement and proof are the article's own text line for line. Required context retained uncounted: none. Every definition, display and label of those lines is retained unchanged. Omitted from this package and not redistributed: 1--1590, 1594--1594, 1649--3583. Nothing inside a retained excerpt is omitted or altered: every retained body is delivered exactly as the article's lines are written, so no omission receipt of any kind accompanies this package. Citations: the retained excerpts carry no citation command at all, so no bibliography entry of the article is transcribed and no reference is redistributed. Reference adaptation: none was needed and none was made; every reference inside the delivered unit resolves inside it, so no unresolved reference or citation is produced. Printed slips kept literal: no printed word, symbol, hypothesis or number of the counted statement or of its proof was altered. Layout and class: the article's own class is \texttt{amsart} with packages including amssymb, amsthm, amsmath, graphicx, esvect, xcolor, tikz, hyperref, amsrefs, enumitem, geometry, standalone, microtype, fontenc; the wrapper is the lane's pinned \texttt{amsart} editorial wrapper at 10pt on A4, which restates the theorem-like environments the retained text needs and adds the article's own macro definitions that the retained text needs (\texttt{\textbackslash arc}), with extra wrapper packages (bm, bbm, dsfont, xspace, cancel, mathtools). The delivered numbering is the unit's own. Assets: no retained span contains a figure environment, a graphics-inclusion command, a PDF-page inclusion, a table or a TikZ picture; no image file is copied into this package, the package declares no asset dependency and no image file is redistributed with it. Licence: version arXiv:2402.12643v2 of this article is distributed under the Creative Commons Attribution 4.0 International licence, as recorded in the arXiv licence metadata for this exact version and shown on the abstract page https://arxiv.org/abs/2402.12643v2. A full rights notice is delivered at licenses/arxiv-2402.12643-CC-BY-4.0.txt. Characters: every retained excerpt is pure ASCII, so the delivered engine, XeLaTeX with its default Latin Modern text font, reports no missing character.
\begin{proposition}\label{piorientation}
Any function $f\colon\partial P\mapsto \partial P$ that maps $\arc[x,f(x)]$ into $\arc[f(x),x]$ for all nonfixed points $x$ of $f$ is orientation preserving. In particular, the Poncelet map is orientation preserving.    
\end{proposition}

\begin{proof}
Let $f\colon\partial P\mapsto \partial P$ be such that it maps $\arc[x,f(x)]$ into $\arc[f(x),x]$ for all $x$ not fixed by $f$.
Observe that the property assumed for $f$ is precisely (ii) of the previous lemma. In the previous lemma, (iii) follows from (ii). We thus have that $f(z)\in\arc[z,f(x)]$ for all $x$ not fixed by $f$ and $z\in \arc(f(x),x]$. 

We will make use of another property of $f$ regarding fixed points: Let $y\in \partial P$ be fixed by $f$.
If $x$ is not fixed by $f$, then $\arc(x,f(x))$ does not contain any fixed points of $f$ (as it is mapped into its complement).
It follows that $y\in \arc[f(x),x]$, i.e., $f(x)\in \arc[x,y]$. This conclusion holds  for all $x\neq y$ (fixed or not fixed by $f$). 


Let $x\neq y$ be such that $f(x)\neq f(y)$, and let $z\in \arc(x,y)$. We wish to show that  $f(z)\in \arc[f(x),f(y)]$. 
We examine several cases.

\textbf{Case that $y$ is fixed by $f$}. As remarked above, in this case  $f(x)\in\arc[x,y]$ and $f(z)\in \arc[z,y]$. We either have that  $z\in \arc(x,f(x))$ or that $z\in [f(x),y)$. If $z\in\arc[x,f(x)]$, then 
\[
f(z)\in \arc[f(x),x]\cap\arc[z,y]=\arc[f(x),f(y)].
\] 
If $z\in \arc [f(x),y]$, then $f(z)\in \arc[z,y]\subseteq \arc[f(x),f(y)$. 

\textbf{Case that $x$ is fixed by $f$}. If $f(z)=f(x)$, then we certainly have that $f(z)\in\arc[f(x),f(y)]$, as desired.
Assume thus that $f(z)\neq f(x)$. Notice that  $y\in \arc(z,x)$ and that $x$ is fixed by $f$. We can thus apply the previous case 
to the triple of points $z,y,x$. We deduce that $f(y)\in \arc[f(z),f(x)]$, i.e., that $f(z)\in \arc[f(x),f(y)]$, as desired.


In the remaining cases we assume that neither $x$ nor $y$ is a fixed point of $f$.


\textbf{Case  $x\in \arc[y,f(y)]$}. In this case $f(x)\in \arc[f(y),y]$. If $z\in \arc[x,f(x)]$, then
\[
f(z)\in \arc[f(x),x]\subseteq \arc[f(x),f(y)].
\]
If on the other hand $z\in \arc(f(x),y]\subseteq \arc(f(y),y]$, then 
\[
f(z)\in \arc[z,f(y)]\subseteq \arc[f(x),f(y)].
\]


\textbf{Case $x\in \arc(f(y),y)$}. 
In this case, $f(x)\in [x,f(y)]$. We  either have that $f(x)\in (x,y]$ or that $f(x)\in (y,f(y))$. We analyze each of these cases next.

\textbf{Subcase $f(x)\in (x,y]$}. If $z\in \arc[x,f(x)]$, then
\[
f(z)\in \arc[f(x),x]\cap \arc[z,f(y)]=\arc[f(x),f(y)],
\] 
and if $z\in \arc (f(x),y]\subseteq (f(y),y]$, then 
\[
f(z)\in \arc[z,f(y)]\subseteq \arc[f(x),f(y)].
\]


\textbf{Subcase $f(x)\in (y,f(y))$}. Then
\[
f(z)\in \arc[f(x),x]\cap \arc[z,f(y)]=\arc[f(x),f(y)].\qedhere
\]
\end{proof}

\end{document}
